\documentclass[12pt]{article}

%% Basic document formatting
\usepackage{amsmath, amsthm, amssymb, amsfonts}
\usepackage{mathtools}
\usepackage{mathrsfs}
\usepackage{xspace}
\usepackage{thmtools}
\usepackage{graphicx}
\usepackage{tikz-cd}
\usepackage{color}
\usepackage{setspace}
\usepackage{fancyhdr}
\usepackage{titling}
\usepackage{hyperref}
\usepackage[left=0.4in,right=0.4in,top=1in,bottom=1in]{geometry}
\usepackage{float}
\usepackage{tabularx}
\usepackage{booktabs}
\usepackage[utf8]{inputenc}
\usepackage[english]{babel}
\usepackage{framed}
\usepackage[dvipsnames]{xcolor}
\usepackage{environ}
\usepackage{tcolorbox}
\tcbuselibrary{theorems,skins,breakable}

\usepackage{enumitem}
\setlist[enumerate]{leftmargin=*}
\setlist[enumerate,1]{labelindent=\parindent}
\setlist[enumerate,2]{labelindent=0pt}

% Blackboard Bold
\newcommand{\Z}{\mathbb{Z}}                 % integers
\newcommand{\Zpl}{\mathbb{Z}_{+}}           % positive integers
\newcommand{\Znn}{\mathbb{Z}_{\geq 0}}      % nonnegative integers. 
\newcommand{\Q}{\mathbb{Q}}                 % rationals
\newcommand{\Qpl}{\mathbb{Q}_{+}}           % positive Rationals
\newcommand{\R}{\mathbb{R}}                 % reals
\newcommand{\Rpl}{\mathbb{R}_{+}}           % positive Reals
\newcommand{\C}{\mathbb{C}}                 % complex numbers
\newcommand{\F}{\mathbb{F}}                 % field

% Words
\newcommand{\st}{\text{ such that }}
\newcommand{\wrt}{\text{ with respect to }}
\newcommand{\with}{\text{ with }}
\newcommand{\ie}{\text{, i.e. }}
\newcommand*{\ditto}{\text{---}\texttt{"}\text{---}}

% Operators
\newcommand{\abs}[1]{\left|#1\right|}                   % absolute value
\newcommand{\norm}[1]{\abs{\abs{#1}}}
\newcommand{\floor}[1]{\left\lfloor #1 \right\rfloor}   % floor
\newcommand{\ceil}[1]{\left\lceil #1 \right\rceil}      % ceiling
\newcommand{\powset}[1]{\mathscr{P}(#1)}

% Category Theory 
\renewcommand{\hom}{\text{Hom}}
\newcommand{\homspace}[3]{\hom_{#1}(#2, #3)}
\newcommand{\coproduct}[9]{
  \begin{tikzcd}[ampersand replacement=\&]
      \& \& #1 \arrow[dll, bend right, "#6"] \arrow[dl, "#5"] \\
   #4 \& #3 \arrow[l, "#9"] \\
      \& \& #2 \arrow[ull, bend left, "#8"] \arrow[ul, "#7"]
  \end{tikzcd}
}

\newcommand{\product}[9]{ 
  \begin{tikzcd}[ampersand replacement=\&]
      \& \& #3 \\
      #1 \arrow[r, "#7"] \arrow[urr, bend left, "#5"] \arrow[drr, bend right, "#6"] \& #2 \arrow[ur, "#8"] \arrow[dr, "#9"] \\ 
      \& \& #4
  \end{tikzcd}
}


% Algebra 
\newcommand{\Syl}[2]{\text{Syl}_{#1}{(#2)}}           % set of Sylow-p subgroups 
\newcommand{\<}{\langle}                % \<x\>, subgroup generated by x 
\renewcommand{\>}{\rangle}
\renewcommand{\index}[2]{[#1 : #2]}
\newcommand{\id}{\text{id}}             % identity element
\newcommand{\lhdneq}{%
  \mathrel{\ooalign{$\lneq$\cr\raise.22ex\hbox{$\lhd$}\cr}}}
\newcommand{\order}[1]{\text{o}(#1)}    % order of an element
\newcommand{\spec}{\operatorname{Spec}}
\newcommand{\rad}{\operatorname{rad}}   % radical of an ideal 
\newcommand{\sym}[1]{\mathfrak{S}_{#1}}
\newcommand{\aut}[1]{\text{Aut}(#1)} 
\newcommand{\gal}[2]{\text{Gal}(#1 / #2)} 
\newcommand{\inn}[1]{\text{Inn}(#1)} 
\let\oldcong\cong
\let\oldequiv\equiv
\renewcommand{\cong}{\oldequiv}
\renewcommand{\equiv}{\oldcong}

\newcommand{\triangleleftneq}{\mathrel{\ooalign{%
  $\trianglelefteq$\cr
  \hidewidth\raise0.06ex\hbox{$\not\mkern4mu$}\hidewidth\cr
}}}

\makeatletter % cycle
\newcommand{\cyc}[1]{(\mathbf{\cyc@process#1\relax})}
\def\cyc@process#1#2\relax{%
	#1%
	\ifx\relax#2\relax
	\else
		\,\cyc@process#2\relax
	\fi
}
\makeatother

\newcommand{\GL}[2]{\text{GL}_{#1}(#2)}                             % general linear group
\newcommand{\SL}[2]{\text{SL}_{#1}(#2)}                             % special linear group
\newcommand{\gl}[2]{\mathfrak{gl}_{#1}(#2)}
\renewcommand{\sl}[2]{\mathfrak{sl}_{#1}(#2)}
\newcommand{\so}[2]{\mathfrak{so}_{#1}(#2)}
\newcommand{\su}[1]{\mathfrak{su}(#1)}

% Linear Algebra
\newcommand{\bmat}[1]{\begin{bmatrix}#1\end{bmatrix}}   % bracketed matrix
\newcommand{\rank}{\operatorname{rank}}                 % rank
\newcommand{\nullity}{\operatorname{nullity}}           % nullity
\newcommand{\tr}{\operatorname{tr}}

% Combinatorics
\newcommand{\qnum}[1]{\left[#1\right]_q}                % q-analog
\newcommand{\qfact}[1]{\left[#1\right]!_q}              % q-analog factorial 
\newcommand{\qbinom}[2]{\binom{#1}{#2}_q}               % Gaussian binomial coefficient

% Topology/Analysis
\newcommand{\ball}[2]{\text{B}_{#1}(#2)}  % B_r(x): open r-balls around x
\newcommand{\diam}{\text{diam}}           % diamter of a set in metric space 
\newcommand{\lowint}[2]{\underline{\int_{#1}^{#2}}}
\newcommand{\upint}[2]{\overline{\int}_{#1}^{#2}}

% Misc Notation
\newcommand{\oneton}[1]{\{1, \ldots, #1\}}
\newcommand{\defeq}{\vcentcolon=}   % :=
\newcommand{\eqdef}{=\vcentcolon}   % =: 
\renewcommand{\bf}[1]{\textbf{#1}}
\newcommand{\im}{\operatorname{im}}
\newcommand{\fnrest}[2]{\left.#1\right|_{#2}}
\newcommand{\isom}{\overset{\sim}{\rightarrow}} 


\renewcommand{\thesection}{\arabic{section}.}
\renewcommand{\thesubsection}{(\alph{subsection})}

\newtheorem{claim}{Claim}
\newtheorem*{lemma}{Lemma}

%% Headers & title setup
\newcommand{\course}{Math100C}
\newcommand{\myname}{Jay Ser}
\setlength{\headheight}{14.5pt}
\pagestyle{fancy}
\fancyhf{}
\renewcommand{\headrulewidth}{0.4pt}
\lhead{\course}
\rhead{\myname}
\cfoot{\thepage}
\setlength{\droptitle}{-4em} 
\title{\course\ - WK5} 
\author{\myname}
\date{2026.05.01}

\begin{document}
\maketitle
\thispagestyle{fancy}

%------------------------------------------------------------------------------%


\section{} 
Let $P(x)$ be a polynomial over $\mathbb{R}$ of \emph{odd} degree. 
Without using the fundamental theorem of algebra, prove that $P(x)$ has a root over $\mathbb{R}$; 
in particular, there is no irreducible polynomial over $\mathbb{R}$ of odd degree greater than $1$.
\emph{(Hint: use the intermediate value theorem.)}

\noindent\dotfill

Without loss of generality, assume the leading coefficient of $P$ is positive. 
Recall from grade school that polynomials are continuous, and for polynomials $P(x)$ of odd degree with a positive leading coefficient,
$$P(x) \rightarrow \infty \text{ as } x \rightarrow \infty, \hspace{0.5em} P(x) \rightarrow -\infty \text{ as } x \rightarrow -\infty.$$
Hence the intermediate value theorem applies to give some $x_0 \in \R$ such that $P(x_0) = 0$,
i.e., the intermediate value theorem guarantees a root of $P$. 
 
\section{} 
An \emph{algebraic integer} is an element of $\mathbb{C}$ which is algebraic over $\mathbb{Q}$ and which is the root of some monic polynomial with \emph{integer} coefficients.
\subsection{}
Prove that $\alpha \in \mathbb{C}$ is an algebraic integer if and only if $1, \alpha, \alpha^2, \dots$ belong to some finitely generated $\mathbb{Z}$-submodule of $\mathbb{C}$.
\emph{(Hint: follow the argument for algebraic numbers given in lecture. 
You will need to use the fact from 100B that $\mathbb{Z}$ is a \emph{noetherian ring} to argue that the submodule generated by the $\alpha^i$ is itself finitely generated.)}

\noindent\dotfill 

% There is a unique minimal irreducible monic polynomial in $\Z[x]$ satisfied by $\alpha$: 
% by definition of an algebraic integer, there is a monic polynomial $f(x) \in \Z[x]$ satisfied by $\alpha$, and because $\Z[x]$ is a unique factorization domain, $f(x)$ factors into monic irreducible factors such that one of them is satisfied by $\alpha$; 
% the irreducible monic polynomial in $\Z[x]$ satisfied by $\alpha$ is unique by the same proof for the existence of \textit{the} minimal monic polynmomial for an algebraic element over a field using the Euclidean algorithm for polynomials. 
Suppose $\alpha$ is an algebraic integer. 
Let $\alpha$ satisfy $f(x) = x^n + a_{n - 1}x^{n - 1} + \ldots + x_0 \in \Z[x]$. 
The identification 
\begin{equation} \label{eq:2_Zmod}
  \Z[x] / \<f(x)\> \equiv \{c_0 + c_1 \alpha + \ldots + c_{n - 1}\alpha^{n - 1} \mid c_{i} \in \Z\}
\end{equation}
makes $\Z[x] / \<f(x)\>$ a finitely generated $\Z$-module. 
Further, the identification \eqref{eq:2_Zmod} shows that $\Z[x] / \<f(x)\>$ contains all the powers of $\alpha$:
Clearly $\Z[x] / \<f(x)\>$ contains $1, \ldots, \alpha^{n - 1}$, and that it contains $\alpha^d$ for $d \geq n$ is given inductively as follows;
the point is that each $\alpha^d$ can be written as a $\Z$-linear combination of $1, \alpha, \ldots, \alpha^{n - 1}$ and hence a member of $\Z[x] / \<f(x)\>$ by \eqref{eq:2_Zmod}. 
For $d = n$, the identity $f(\alpha) = 0$ means 
\begin{equation} \label{eq:2_alphapower}
  \alpha^n = -(a_{n - 1}\alpha^{n - 1} + \ldots + a_0)
\end{equation}
in $\Z[x] / \<f(x)\>$, which means $\alpha^n$ is indeed an element of the $\Z$-module identified in the $\Z$-module on the right side of \eqref{eq:2_Zmod}. 
Next, if $\alpha^{d - 1}$ can be written as a $\Z$-linear combination of $1, \alpha, \ldots, \alpha^{n - 1}$, say with coeffients $c_i \in \Z$, then writing 
$$\alpha^{d} = \alpha^{d - 1}\alpha = (c_0 + c_1 \alpha + \ldots + c_{n - 1}\alpha^{n - 1})\alpha$$
and using the simplification for $\alpha^n$ as given in \eqref{eq:2_alphapower} shows that $\alpha^d$, too, can be written as a $\Z$-linear combination of $1, \alpha, \ldots, \alpha^{n - 1}$, and hence a member of $\Z[x]/\<f(x)\>$. 
This proves the forward direction. 

Conversely, suppose a finitely generated $\Z$-submodule $M$ of $\C$ contains all the powers $1, \alpha, \alpha^2, \ldots$.  
Since $\Z$ is a Noetherian ring and $M$ is finitely generated, $M$ is a Noetherian $\Z$-module. 
Hence the submodule of $M$ generated by the power of $\alpha$ -- call this submodule $L$ -- is finitely generated. 
That $L$ is finitely generated means that there is an integer $n$ such that $\alpha^{n}$ is $\Z$-linearly dependent to $1, \alpha, \ldots, \alpha^{n - 1}$; 
in other words, there exists $c_0, \ldots, c_{n - 1}$ such that 
$$\alpha^{n} - c_{n - 1} \alpha^{n - 1} - \ldots - c_0 = 0.$$
This shows that $\alpha$ is an algebraic integer. 

\subsection{}
Deduce that the algebraic integers form a subring of $\mathbb{C}$.
 
\noindent\dotfill 

Let $\alpha$ and $\beta$ be algebraic integers, $\alpha$ of degree $d_1$ and $\beta$ of degree $d_2$.  
There is a finitely generated $\Z$-submodule of $\C$ that contains all the powers of $\alpha$ and $\beta$, say $L$.  
Since $\Z$ is a Noetherian ring, $L$ is a Noetherian module. 
Accordingly, the collection 
$$\{\alpha^i \beta^j\}_{\substack{i \in \Znn \\ j \in \Znn}}$$ 
generates a finitely generated $\Z$-submodule of $L$. 
This submodule of $L$ contains all the powers of $\alpha - \beta$ and $\alpha \beta$, so $\alpha - \beta$ and $\alpha \beta$ are algeraic integers. 
This shows that the algebraic integers form a ring. 


\section{} 
Use Gauss's lemma from 100B to prove that $\alpha \in \mathbb{C}$ is an algebraic integer if and only if $\alpha$ is an algebraic number and its minimal polynomial has integer coefficients. In particular, a rational number is an algebraic integer if and only if it is an integer. (This will be used in the optional exercises about character tables.)
 
\noindent\dotfill 

The forward direction holds immedaitely. 
Let $\alpha$ be an algebraic number whose minimal polynomial is $f(x) = x^n + + a_{n - 1} x^{n - 1} \ldots + a_0 \in \Q[x]$ where each $a_i$ is an integer. 
By definition, $f$ is monic and irreducible in $\Q[x]$. 
Gauss's Lemma says a monic polynomial with integer coefficients is irreducible in $\Q[x]$ if and only if it is irreducible in $\Z[x]$, so $f$ is irreducible in $\Z[x]$ as well. 
Hence $\alpha$ is an algebraic integer satisfying a monic irreducible polynomial in $\Q[x]$. 

The second statement follows from the ratioanl root theorem: 
if $\frac{s}{t}$ is a rational root of an integer polynomial $a_{n}x^{n} + \ldots + a_0$, then $t \mid a_n$ and $s \mid a_0$. 
Hence if $\frac{s}{t}$ is algebraic integer, then $t \mid 1$ because $\frac{s}{t}$, by definition, satisfies a monic polynomial. 

\section{} 
Compute the minimal polynomials over $\mathbb{Q}$ of the following complex numbers:
\[
  i + \sqrt{2}, \qquad i + \sqrt{3}, \qquad \sqrt{2} + \sqrt{3}.
\]
Be sure to check that your polynomials are irreducible! \emph{(Hint for that: if a polynomial of degree $4$ over $\mathbb{Q}$ factors into two quadratic factors, then the sum \emph{and} product of the two roots of either factor must be rational numbers.)}

\noindent\dotfill 

I first prove Kiran's hint. 
A posteriori, knowing how the polynomials in the following solutions are found, one need not the result. 
However, seeing the polynomials in a vacuum, one needs some methods to determine whether or not they are irreducible. 
\begin{lemma}
  The sum and product of the roots of a rational quadratic polynomial  both lie in $\Q$. 
\end{lemma}
\begin{proof}[Proof of Lemma]
  Take $f(x) = ax^2 + bx + c \in \Q[x]$. 
  The quadratic formula says that the two roots of $f(x)$ are given by 
  $$\frac{-b \pm \sqrt{b^2 - 4ac}}{2a}.$$
  Letting $\alpha_1$ and $\alpha_2$ denote the roots of $f(x)$, 
  $$\alpha_1 + \alpha_2 = -\frac{b}{a}, \hspace{1.5em} \alpha_1 \alpha_2 = \frac{1}{4a^2}(b^2 - (b^2 - 4ac)) = \frac{c}{a}$$
  so $\alpha_1 + \alpha_2, \alpha_1\alpha_2 \in \Q$. 
\end{proof}

\subsection{}
Denote $\alpha = i + \sqrt{2}$. 
One can check 
$$(x + i + \sqrt{2})(x + i - \sqrt{2})(x - i + \sqrt{2})(x - i - \sqrt{2}) = x^4 - 2x^2 + 9,$$
so $f(x) \defeq x^4 - 2x^2 + 9$ has $\alpha$ as a root. 
By the rational root theorem, if $f$ has a root, then it is $\pm 1, \pm 3$, or $\pm 9$, and clearly none of them satisfy $f$. 
This shows $f$ doesn't have linear and degree 3 factors. 
By the lemma, $f$ also doesn't have quadratic factors because no pair from $\pm i \pm \sqrt{2}$ simultaneously add and multiply to an integer (and the hypothetical quadratic factor MUST have one of these as a factor since $\C[x]$ is a UFD): 
the roots sum to an integer if and only if they sum to 0, which occurs if and only if one root is the negative of the other, but then this pair doesn't multiply to an integer.
Hence $f$ is irreducible. 

\subsection{} 
Denote $\alpha = i + \sqrt{3}$. 
One can check 
$$(x + i + \sqrt{3})(x + i - \sqrt{3})(x - i + \sqrt{3})(x - i - \sqrt{3}) = x^4 - 4x^2 + 16$$
so $f(x) \defeq x^4 - 4x^2 + 16$ has $\alpha$ as a root. 
By the rational root theorem, the only possible rational roots for $f$ are $\pm 1, \pm 2, \pm 4, \pm 8, \pm 16$. 
Simple calculations show that $\pm 1$ and $\pm 2$ don't satisfy $f$, and clearly the bigger integers don't satisfy $f$ either because the $x^4$ balloons everything. 
And $f$ doesn't have a quadratic factor because no pair from $\pm i \pm \sqrt{3}$ simultaneously add and multiply to an integer. 
Hence $f$ is irreducible. 

\subsection{}
Denote $\alpha = \sqrt{2} + \sqrt{3}$. 
One can check 
$$(x - \sqrt{2} - \sqrt{3})(x - \sqrt{2} + \sqrt{3})(x + \sqrt{2} - \sqrt{3})(x + \sqrt{2} + \sqrt{3}) = x^4 - 10x^2 + 1,$$
so $f(x) \defeq x^4 - 10x^2 + 1$ has $\alpha$ as a root. 
By the rational root theorem, the only possible rational roots for $f$ are $\pm 1$ and clearly they aren't roots of $f$. 
If $f$ is reducible, then it can only have an irreducible quadratic polynomial as a factor. 
But no pair from $\{\pm \sqrt{2} \pm \sqrt{3}\}$ both add and multiply to an integer, so the lemma says $f$ can't have any quadratic factors either. 
$f$ is irreducible. 
 
\section{}
 
Let $p$ be a prime number. Show that the polynomial $((x+1)^p - 1)/x$ is irreducible over $\mathbb{Q}$ using Gauss's lemma and the Eisenstein irreducibility criterion from 100B. Then deduce that $e^{2\pi i/p}$ is an algebraic number over $\mathbb{Q}$ of degree $p-1$. (Note that this works even for $p=2$!) \emph{Optional: do the analogue for $e^{2\pi i/p^n}$ for $n > 1$.}

\noindent\dotfill 

Computation shows 
$$((x + 1)^{p} - 1)/x = \sum_{k = 1}^{p} \binom{p}{k} x^{k - 1}, $$
and that $p \mid \binom{p}{k}$ for any $k \in \{1, \ldots, p - 1\}$ shows that the Eisenstein criterion applies; 
$((x + 1)^{p} - 1)/x$ is irreducible over $\Q$. 

Let $\zeta_{p} = e^{e\pi i / p}$. 
To prove that $\zeta_{p}$ is an algebraic number of degree $p - 1$, it suffices to show that $\Phi(x) \defeq x^{p - 1} + x^{p - 2} + \ldots + 1$ is the minimal polynomial for $\zeta_p$ over $\Q$. 
One sees $\Phi(\zeta_p) = 0$ by noting 
\begin{equation} \label{eq:5_factor}
  \Phi(x) (x - 1) = x^p - 1;
\end{equation}
since $\zeta_p$ is a root of the polynomial on the right hand side but isn't a root of $x - 1$, $\Phi(\zeta_p) = 0$ is necessary.   
It remains to show that $\Phi$ is irreducible. 
Notice that substituting $x + 1$ for $x$ in \eqref{eq:5_factor} yields
$$(x + 1)^p - 1 = x \Phi(x + 1),$$
or dividing both sides $x$, 
$$\Phi(x + 1) = ((x + 1)^{p} - 1) / x.$$
The previous paragraph showed that this exact polynomial is irreducible.
But because the algebraic substitution 
$$\alpha: \Q[x] \mapsto \Q[x]; \hspace{0.5em} \fnrest{\alpha}{\Q} = \id, \, x \mapsto x + 1$$ 
is an automorphism on $\Q[x]$ (for example, the substitution map $x \mapsto x - 1$ is its inverse), 
the irreducibility of $\Phi(x + 1)$ is equivalent to the irreducibility of $\Phi(x)$. 

This completes the proof that $\zeta_p$ is an algebraic number of degree $p - 1$; 
particularly, $\Phi(x) = x^{p - 1} + x^{p - 2} + \ldots + 1$ is its minimal polynomial over $\Q$. 
 
\section{}
 
Show that every rational number appearing in the character table of a finite group is actually an integer. 
Then use the week 4 optional homework to show that the character table of $\mathfrak{S}_n$ consists entirely of integers.

\noindent\dotfill 

Let $\chi$ be a character of a finite group $G$. 
Recall that for any $g \in G$, $\chi(g)$ is the sum of roots of unity. 
The roots of unity are algebraic integers, and since by Q2 the algebraic integers form a ring, $\chi(g)$ is an algebraic integer as well. 
Q3 says rational algebraic integers are integers. 
Hence rational entries in a character table are integers. 

In $\mathfrak{S}_n$, two elements generate the same subgroup only if they are of the same cycle type; 
this means two elements generate the same subgroup only if they are in the same conjugacy class. 
By HW4 Q4, the entries for the character table of $\mathfrak{S}_n$ must be integers. 
 
\section{}
 
Let $G$ be a finite group. Let $C_1, \dots, C_n$ be the conjugacy classes of $G$. Fix an irreducible character $\chi$ of $G$ and set
\[
  \omega_i := \sum_{g \in C_i} \frac{\chi(g)}{\chi(1)} \qquad (i = 1, \dots, n).
\]
Prove that the set of $\mathbb{Z}$-linear combinations of $\omega_1, \dots, \omega_n$ is a subring of $\mathbb{C}$;
then deduce from this that the $\omega_i$ are algebraic integers.

\noindent\dotfill 

Let $Z$ be the set of all $\Z$-linear combinations of $\omega_1, \ldots, \omega_n$. 
Let $R$ be the representation associated with $\chi$. 
The $\omega_i$'s are closed under addition by definition. 
The following shows that the $\omega_i$'s are closed under multiplication as well.
Take any $\omega_i$ and $\omega_j$. 
The calculation 
\begin{align*}
  (\sum_{g \in C_i} R_g)R_h & = \sum_{g \in C_i} R_{gh} \\ 
                            & = \sum_{g \in C_i} R_{h h^{-1}gh} \\ 
                            & = R_h \sum_{g \in C_i} R_g 
\end{align*}
shows that $\sum_{g \in C_i} R_g$ is a $G$-invariant linear transformation. 
Since $\chi$ is irreducible, Schur's lemma says 
\begin{equation} \label{eq:7_sumRscalar}
  \sum_{g \in C_i} R_g = \lambda_i I_n
\end{equation}

for some $\lambda_i \in \C$. 
It follows that 
\begin{equation} \label{eq:7_omegaval}
  \omega_i = \frac{1}{\chi(1)} \text{tr} \sum_{g \in C_i} R_g = \frac{1}{\chi(1)}\lambda_i\chi(1) = \lambda_i
\end{equation}

Take some $g \in C_i$ and some $h \in C_j$. 
Let $C_{gh}$ be the conjugacy class of $gh$. 
Then for any $x \in G$, 
$$xghx^{-1} = xgx^{-1} x hx^{-1}$$ 
shows that $xghx^{-1}$ is the product of an element in $C_i$ and an element in $C_j$. 
Hence summing over all $g \in C_i$ and $h \in C_{j}$, 
\begin{equation} \label{eq:7_sum}
  \sum_{g \in C_i} \sum_{h \in C_j} R_g R_h = \sum_{g \in C_i} \sum_{h \in C_j} R_{gh} = \sum_{k = 1}^{n} m_k \sum_{g \in C_k} R_g,
\end{equation}
whether the $m_k$'s count the number of different combinations of $g \in C_i$ and $h \in C_j$ such that $gh$ lands in the conjugacy class $C_k$. 
Particularly, $m_k$ is an integer. 
Because \eqref{eq:7_sumRscalar} shows that $\sum_{g \in C_i} R_g$, etc. are scalar matrices, one can calculate trace as follows: 
$$\text{tr}(\sum_{g \in C_i} R_g \sum_{h \in C_j} R_h) = \lambda_i \lambda_j \chi(1) = \omega_i \omega_j \chi(1).$$
Combining this with \eqref{eq:7_sum} and \eqref{eq:7_omegaval},
$$\omega_i \omega_j \chi(1) = \text{tr} \sum_{k = 1}^{n} m_k \sum_{g \in C_k} R_g = \sum_{k = 1}^{n} m_k \omega_k \chi(1)$$
Cancelling out the $\chi(1)$ in both side shows that $\omega_i \omega_j$ is a $\Z$-linear combination of the $\omega_k$'s, hence $Z$ is closed under multiplication; 
$Z$ is a ring. 

Namely, that $Z$ is a ring mean that all the powers of each of the $\omega_i$'s lie in $Z$. 
$Z$ is clearly a finitely generated $\Z$-module, hence the $\omega_i$'s are algebraic integers. 

\section{}
 
With notation as in the previous exercise, choose any $g_i \in C_i$. Prove that
\[
  |G| = \sum_{i=1}^n \chi(1)\, \omega_i\, \chi(g_i^{-1}),
\]
then deduce the missing assertion from the orthogonality theorem: the dimension of every irreducible representation of $G$ divides the order of $G$.

\noindent \dotfill 

Notice that for each $i \in \oneton{n}$, 
$$\omega_i = \sum_{g \in C_i} \frac{\chi(g)}{\chi(1)} = \#C_i \frac{\chi(g_i)}{\chi(1)}$$
Also, recall $\chi(h^{-1}) = \overline{\chi(h)}$ for any $h \in G$ and any character $\chi$.  
Hence 
\begin{align*}
  \sum_{i = 1}^{n} \chi(1) \omega_i \chi(g_{i}^{-1}) 
  & = \sum_{i = 1}^{n} \#C_i \chi(g_i) \overline{\chi(g_i)} \\ 
  & = \#G \<\chi, \chi\>  = \#G 
\end{align*}
where the $\<\chi, \chi\> = 1$ because $\chi$ is irreducible. 
This verifies the desired formula; 
namely, $\#G$ is a multiple of $\chi(1)$, i.e., the dimension of every irreducible character divides the order of the group. 

\end{document}
